\section{Related Work}
\label{sec:related}

\subsection{HIV bnAb measurement and development}
CATNAP standardizes antibody-by-virus neutralization panels and associated
metadata~\cite{yoon2015catnap}, curating IC50 and IC80 measurements from
pseudovirus neutralization assays across antibody-virus pairs and
providing the consistent numerical labels for our regression target.
Longitudinal lineage studies show that breadth can change substantially as
somatic variants accumulate, including within the CAP256/VRC26
lineage~\cite{doriarose2014developmental,bhiman2015cap256}, where members share
high sequence identity yet span a wide range of neutralization phenotypes.
These observations motivate lineage-held-out evaluation: high sequence
similarity does not imply identical breadth, and lineage membership is itself a
strong shortcut that a naive random split that a model could exploit without
learning generalizable features.

\subsection{Machine learning for HIV neutralization}
Most computational work on HIV bnAbs predicts \emph{virus} resistance to a fixed
antibody panel rather than scoring antibody breadth from sequence; this
review surveys this literature~\cite{danaila2022review}. In this framing the model
input is typically viral envelope[is viral envelope explained?]
features and the output is a binary or
continuous resistance label for a specific antibody.
Not as useful for ranking which antibody to deploy for diverse
populations. Some recent methods sharpen the virus centered view.
Multi-task learning across antibodies mitigates HIV-1 subtype bias and captures
co-resistance patterns shared across the antibody panel~\cite{igiraneza2024multitask},
and BRAVE[was the BRAVE acronym written?]
couples ESM2 embeddings of both antibody and virus sequences with a
random-forest classifier to predict resistance across 33
bnAbs~\cite{elanbari2025brave}. We evaluate the converse, antibody centered
task, if ESM3 embeddings provide a richer representation
antibody neutralization breadth, from their heavy-chain sequences.


\subsection{Protein and antibody language models}
ESM2 showed large-scale protein representation learning via masked
language modeling on evolutionary-scale sequence corpora, yielding embeddings
competitive with structure prediction benchmarks~\cite{lin2023esm2}. ESM3 extended
this approach across protein sequence, structure, and function features.
through generative, transformer, with a robust masking algorithm.
The novel feature modeling enabled joint reasoning over
residue, 3D coordinates, and functional annotations~\cite{hayes2025esm3}.
In contrast, antibody models use sequence corpora tailored to
the unusual statistics of the immunoglobulin variable regions.
AntiBERTy was trained it on hundreds of millions of natural antibody sequences,
this and showed the representation space can cover maturation
trajectories~\cite{ruffolo2021antiberty}.[was this shhown?]
AbLang trained separate models
for heavy and light chains using OAS, reflecting the distinct mutational
statistics of each chain, and demonstrates an antibody-specific advantage for
sequence restoration~\cite{olsen2022ablang}. More recent antibody-specific
models scale further and target known biases: IgBert and IgT5 use
over two billion unpaired and two million paired OAS
sequences~\cite{kenlay2024igbert}, while AbLang-2 targets
germline bias, where antibody models tend to be
dominated by proximity to the germline sequence rather than functional
features that might not generalize to downstream tasks~\cite{olsen2024ablang2}.
Observed Antibody Space provides clean and annotated repertoire sequences
used in pre-training and often tuning~\cite{olsen2022oas}.

\subsection{General versus domain-specific representations}
Whether antibody-specific pretraining is needed is unclear. Across
diverse immune tasks, general protein language models can match or exceed
domain-specific models when comparisons control for model size, training corpus
scale, and layer selection~\cite{deutschmann2024domain}.
This result is non-obvious because
antibody sequences violate the distribution of general protein data,
they are defined by the somatic hypermutation and V(D)J
recombinations that produce sequences rare or absent from other sets.
Large general models may compensate through rich, learned representation.
We prop this application of general protein models such as ESM3 against
antibody feature extracted from the sequences.
We compare
adaptation on the functional task of neutralization breadth, rather than a structural 
task tested in prior surveys.

\subsection{Adaptation and latent interpolation}
Low-Rank Adaptation (LoRA) allows tuning of a frozen model through trainable low-rank
decompositions of the attention weights.
This freezes over 99\% of parameters allowing 
some adaptation a new/shifted distribution~\cite{hu2022lora}. 
We use it on the masked self-supervised decoding task to 
tune ESM3 on 
OAS antibody sequences. We do not use LoRA for tuning on bnAb sampples,
thus, theoretically ESM3 is tuned for antibody sequence.
Variational
autoencoders (VAEs)~\cite{kingma2014vae} learn a regularized Gaussian 
latent space via a KL-divergence penalty on the autoencoder.
This learned latent representation with regularization
theoretically allows linear paths in the embedded space to correspond to 
semantically smooth interpolations between data points, creating the manifold.
However, whether this holds for a given modality depends on the representation quality,
decoder fidelity, and of course, the quality and amount of data.
Our VAE experiment compares latent interpolation with a straight-line in the
VAE latent space, with a similar interpolation in the 
ESM3 embedding space.